
We now combine the results of the previous three sections to evaluate
$I(\phi)$.
The accuracy of our BCLT expansion of $\expect{\post}{\phi(\theta, \xvec_s)}$ will
depend on the following residual function.
%
%%%%%%%%%%%%%%%%%%%%%%%%%%%%%%%%%%%%%%%%%
\begin{defn}\deflabel{bclt_resid}
%
Let $\infoev > 0$ denote the minimum eigenvalue of $\info$, and let $\delta > 0$
be an arbitrarily small number. For a function $\phi(\theta, \xvec_s)$ that
is third-order BCLT-okay, define the residual function
%
\begin{align*}
%
\resid{0}(\xvec_s) :={}&
    1 +
    \norm{\phi(\thetahat, \xvec_s)}_2 +
    \norm{\phigrad{1}(\thetahat, \xvec_s)}_2 +
    \norm{\phigrad{2}(\thetahat, \xvec_s)}_2
\\
%
\resid{1}(\xvec_s) :={}&
    \expect{\normresid{\infoev / 2}{R_1}(\tau)}{
        \norm{\tresid{3}{\phi(\cdot, \xvec_s), \theta, \thetahat}}_2}
\\
%
\resid{2}(\xvec_s) :={}&
    \expect{\normresid{\infoev / 2}{R_1 \bigcup R_2}(\tau)}{\norm{\phi(\theta, \xvec_s)}_2}
\\
%
\resid{3}(\xvec_s) :={}& \expect{\prior(\theta)}{\norm{\phi(\theta, \xvec_s)}_2}
%
\\
\resid{\phi}(\xvec_s) :={}&
    \resid{0}(\xvec_s) + \resid{1}(\xvec_s) + \resid{2}(\xvec_s) + \resid{3}(\xvec_s).
%
\end{align*}
%
\end{defn}
%%%%%%%%%%%%%%%%%%%%%%%%%%%%%%%%%%%%%%%%%
%%%%%%%%%%%%%%%%%%%%%%%%%%%%%%%%%%%%%%%%%



%%%%%%%%%%%%%%%%%%%%%%%%%%%%%%%%%%%%%%%%%
%%%%%%%%%%%%%%%%%%%%%%%%%%%%%%%%%%%%%%%%%
%%%%%%%%%%%%%%%%%%%%%%%%%%%%%%%%%%%%%%%%%
%
\begin{lem}\lemlabel{bayes_clt}
%
Let \assuref{bayes_clt} hold, and assume that $\phi$ is third-order BCLT-okay.
With $\fdist$-probability approaching one,
%
\begin{align*}
%
\MoveEqLeft
\expect{\post}{\phi(\theta, \xvec_s)} - \phi(\thetahat, \xvec_s) =
\\&
N^{-1} \left(
    \frac{1}{2} \phigrad{2}(\thetahat, \xvec_s) \expect{\normhat}{\tau^2} +
    \frac{1}{6} \likhatk{3}(\thetahat) \phigrad{1}(\thetahat, \xvec_s)
        \expect{\normhat}{\tau^4} \right) +
\\&
\ordlog{N^{-2}} \resid{\phi}(\xvec_s),
%
\end{align*}
%
where $\resid{\phi}(\xvec_s)$ is given in \defref{bclt_resid}.
%
\begin{proof}
%
First, combine the results of \lemref{region_r3, region_r2, region_r1_final_bound},
consolidating higher orders into lower orders and expanding the domain of
integration in the residual $\resid{\phi}(\xvec_s)$.
%
\begin{align}
%
I(\phi) ={}&
    \phi(\thetahat, \xvec_s)
    \int_{R_1} \exp\left(
        N \likhat(\theta) - N \likhat(\thetahat)
    \right) d\tau +
\nonumber\\ &\quad
    N^{-1} \normalizer \left(
        \frac{1}{2} \phigrad{2}(\thetahat, \xvec_s) \expect{\normhat}{\tau^2} +
        \frac{1}{6} \likhatk{3}(\thetahat) \phigrad{1}(\thetahat, \xvec_s)
            \expect{\normhat}{\tau^4} \right) +
\nonumber\\ &\quad
    \ordlog{N^{-2}} \resid{\phi}(\xvec_s).
\eqlabel{i_phi_expansion}
%
\end{align}
%
For convenience in the final bound of \thmref{bayes_clt_main} we have included the
term $\norm{\phi(\thetahat, \xvec_s)}_2$ in the definition of the residual even
though it does not occur in \lemref{region_r1_final_bound}.

The final expectation is given by the ratio
%
\begin{align*}
\expect{\post}{\phi(\theta, \xvec_s)} ={}& \frac{I(\phi)}{I(1)}.
\end{align*}
%
For the remainder of this proof, for compactness of notation, we define
%
\begin{align*}
%
\zeta :={}&
    \int_{R_1} \exp\left(N \likhat(\theta) - N \likhat(\thetahat) \right) d\tau \\
={}& \normalizer + \ordlog{N^{-1}} \\
={}& \normalizer (1 + \ordlog{N^{-1}}),
%
\end{align*}
%
which follows from \lemref{r1_normalizer_expansion} and the fact that
$\normalizer$ is bounded away from zero when $\regevent$ occurs. When
$\phi(\theta, \xvec_s) \equiv 1$, then $\resid{0}(\xvec_s) = \resid{1}(\xvec_s) = 0$,
so by \lemref{region_r3, region_r2},
%
\begin{align*}
%
I(1) = \zeta\left(1 + \ord{N^{-3}}\right).
%
\end{align*}
%
We now use the expansion, valid for $|z| < 1$ as $z \rightarrow 0$,
%
\begin{align}\eqlabel{inverse_expansion}
    \frac{1}{1 + z} = 1 - z + O(z^2).
\end{align}
%
Using \eqref{i_phi_expansion} for $I(\phi)$ and
applying \eqref{inverse_expansion} to $1 / I(1)$ and to $1 / \zeta$
gives
%
\begin{align*}
%
\frac{I(\phi)}{I(1)} ={}&
\frac{I(\phi)}{\zeta( 1 + \ord{N^{-3}})} \\
={}& \frac{I(\phi)}{\zeta} + \ord{N^{-3}} I(\phi) \\
={}&
\phi(\thetahat, \xvec_s) +
N^{-1} \frac{\normalizer}{\zeta} \left(
    \frac{1}{2} \phigrad{2}(\thetahat, \xvec_s) \expect{\normhat}{\tau^2} +
    \frac{1}{6} \likhatk{3}(\thetahat) \phigrad{1}(\thetahat, \xvec_s)
        \expect{\normhat}{\tau^4} \right) +
\\&
\ordlog{N^{-2}} \resid{\phi}(\xvec_s)
\\={}&
\phi(\thetahat, \xvec_s) +
N^{-1} \left(
    \frac{1}{2} \phigrad{2}(\thetahat, \xvec_s) \expect{\normhat}{\tau^2} +
    \frac{1}{6} \likhatk{3}(\thetahat) \phigrad{1}(\thetahat, \xvec_s)
        \expect{\normhat}{\tau^4} \right) +
\\&
\ordlog{N^{-2}} \resid{\phi}(\xvec_s).
%
\end{align*}
%
Note that we have absorbed terms from the expansion of $I(\phi)$ that are linear
in $|\phi(\thetahat, \xvec_s)|$, $\norm{\phigrad{2}(\thetahat, \xvec_s)}$, and
$\norm{\phigrad{3}(\thetahat, \xvec_s)}$ into the residual $\resid{\phi}(\xvec_s)$.
To do so we have used the fact that, when $\regevent$ occurs,
$\likhatk{3}(\thetahat) = \likhatk{3}(\thetatrue) + O(1)$.

Finally, the expansion holds with $\fdist$-probability approaching one by
\lemref{regular_event}.
%
\end{proof}
%
\end{lem}
%
%%%%%%%%%%%%%%%%%%%%%%%%%%%%%%%%%%%%%%%%%
